Se presenta el problema de búsqueda de rutas sobre grillas, al cual se incorporan obstáculos y riesgos. Un rover debe atravesar una cuadrícula de N x M desde la esquina superior izquierda hasta la inferior derecha, avanzando únicamente hacia la derecha o abajo, evitando murallas y administrando un número limitado de vidas que se consumen al detonar minas, cada una con una penalizacióñ de P puntos, con el fin de maximizar el puntaje acumulado. Para resolverlo se consideran tres enfoques: la fuerza bruta, que enumera exhaustivamente los $\binom{N + M - 2}{N - 1}$ caminos posible con costo exponencial; la programación dinámica, que aprovecha la subestructura óptima incorporando las vidas restantes al estado y resuelve el problema en $O(NMV)$; y un algoritmo greedy de decisión local, de costo lineal pero in garantía de optimalidad ni de alcanzar la meta.
El objetivo es contrastar el tiempo de ejecución y el consumo de memoria medidos con las cotas asintóticas de cada paradigma, y cuantificar los resultados de greedy mediante su brecha respecto al óptimo y su talla de falla, comparando según tamaño, tipo, y dominio de la grilla
